% =====================================================================
\chapter{Composición e inversa}
\label{tema:4}
% =====================================================================

Componer funciones es aplicarlas una tras otra; invertir una función es deshacerla. Este tema estudia ambas operaciones y cómo se comportan respecto de la inyectividad y la sobreyectividad. Sigue las secciones~12.4 y~12.5 de Hammack~\cite{hammack}, las secciones~5.2 y~5.3 de Velleman~\cite{velleman} y la sección~4.3 de~\cite{mcs}.

\section{Composición}

\begin{definicion}[Composición]
Sean $f\colon A\to B$ y $g\colon B\to C$. La \emph{composición} de $g$ con $f$ es la función $g\circ f\colon A\to C$ definida por $(g\circ f)(a)=g(f(a))$.
\end{definicion}

El orden de lectura es de derecha a izquierda: en $g\circ f$ se aplica primero $f$. La composición solo está definida cuando el codominio de $f$ está contenido en el dominio de $g$.

\addtocounter{definicion}{1}%

\begin{proposicion}\label{prop:t4-asoc}
Sean $f\colon A\to B$, $g\colon B\to C$ y $h\colon C\to D$. Entonces $h\circ(g\circ f)=(h\circ g)\circ f$, y $f\circ\id_A=f=\id_B\circ f$.
\end{proposicion}

\section{Composición, inyectividad y sobreyectividad}

\begin{teorema}\label{teo:t4-comp}
Sean $f\colon A\to B$ y $g\colon B\to C$.
\begin{enumerate}
  \item Si $f$ y $g$ son inyectivas, $g\circ f$ es inyectiva.
  \item Si $f$ y $g$ son sobreyectivas, $g\circ f$ es sobreyectiva.
  \item Si $f$ y $g$ son biyectivas, $g\circ f$ es biyectiva.
\end{enumerate}
\end{teorema}

\begin{proposicion}\label{prop:t4-reciproco}
Sean $f\colon A\to B$ y $g\colon B\to C$.
\begin{enumerate}
  \item Si $g\circ f$ es inyectiva, $f$ es inyectiva.
  \item Si $g\circ f$ es sobreyectiva, $g$ es sobreyectiva.
\end{enumerate}
\end{proposicion}

Las otras dos implicaciones son falsas: si $g\circ f$ es inyectiva, $g$ no tiene por qué serlo, y si $g\circ f$ es sobreyectiva, $f$ no tiene por qué serlo (\cref{eje:4-cambridge}).

\section{Función inversa}

\begin{definicion}[Inversa]
Sea $f\colon A\to B$. Una función $g\colon B\to A$ es \emph{inversa} de $f$ si $g\circ f=\id_A$ y $f\circ g=\id_B$. Si solo se cumple $g\circ f=\id_A$, se dice que $g$ es una \emph{inversa por la izquierda} de $f$; si solo se cumple $f\circ g=\id_B$, una \emph{inversa por la derecha}.
\end{definicion}

\begin{teorema}\label{teo:t4-inversa}
Una función $f\colon A\to B$ tiene inversa si y solo si es biyectiva. En ese caso la inversa es única; se denota $f^{-1}$ y está dada por $f^{-1}(b)=$ el único $a\in A$ con $f(a)=b$.
\end{teorema}

\begin{proposicion}\label{prop:t4-laterales}
Sea $f\colon A\to B$ con $A\neq\emptyset$. Entonces $f$ tiene una inversa por la izquierda si y solo si es inyectiva, y tiene una inversa por la derecha si y solo si es sobreyectiva.
\end{proposicion}

\begin{proposicion}\label{prop:t4-inv-comp}
Si $f\colon A\to B$ es biyectiva, $f^{-1}$ es biyectiva y $(f^{-1})^{-1}=f$. Si además $g\colon B\to C$ es biyectiva, $(g\circ f)^{-1}=f^{-1}\circ g^{-1}$.
\end{proposicion}

\section{Cálculo de inversas}

Para invertir una biyección real $y=f(x)$ se despeja $x$ en función de $y$ y se comprueba que el resultado pertenece al dominio. Cuando $f$ no es biyectiva, puede \emph{restringirse} el dominio y el codominio para obtener una biyección: así se definen $\sqrt{\phantom{x}}$ como inversa de $x^{2}$ en $[0,\infty)$, $\ln$ como inversa de $\exp\colon\R\to(0,\infty)$ y $\operatorname{arcsen}$ como inversa de $\operatorname{sen}$ restringido a $[-\pi/2,\pi/2]$.

\addtocounter{definicion}{1}%

\addtocounter{definicion}{1}%

\addtocounter{definicion}{1}%

